\section{Discussion}
\label{sec:discussion}

\para{What do the regime differences mean?} The simplest interpretation is that the model and the crowd rely on different sources of signal. In \longterm, where only about one crowd prediction is visible on average and more than 100 days remain before resolution, broad world knowledge and textual reasoning may be more useful than crowd momentum. In \shortterm, the key signal is often not in the question text itself but in the sequence of recent updates that encode late-breaking evidence and collective reaction. That is exactly the signal the \questiononly model lacks and the \withhistory model recovers.

\para{What does this imply for LLM deployment?} Our results argue against treating LLMs as uniform replacements for human forecasting crowds. A text-only LLM may be useful as an early-stage independent forecast or as a complement when crowd information is sparse. By contrast, late-stage deployment appears to require structured external signal. In this setup, crowd trajectory summaries provide that signal, but retrieval at the historical cutoff could play a similar role in future work. The practical lesson is that system design should depend on information regime, not on a single aggregate claim about model strength.

\para{How does this relate to prior work?} The results are consistent with the literature review. They align with Halawi \etal's finding that language models can be competitive when the crowd is uncertain or when forecasts are made earlier \cite{halawi2024approaching}, and they remain compatible with Schoenegger and Park's negative result for direct human comparison in realistic forecasting settings \cite{schoenegger2023llm}. Our study adds a controlled within-dataset lens: both patterns can be true at once, depending on horizon and information density.

\para{Limitations.} The study is small, with only 30 items per regime and prompt condition, and none of the paired differences are statistically significant at $\alpha=0.05$. We do not use live retrieval, so the LLM only sees local dataset context rather than the open-web evidence that a real forecaster might consult. The \withhistory condition is not a pure independent forecast because it exposes human trajectory information directly to the model. The \longterm regime is also an especially strict operationalization of low-data, with a mean of only one observed crowd prediction. Finally, the \shortterm sample is dominated by Metaculus questions, which may limit cross-platform generalization.

\para{Broader implications.} Forecasting papers often compress results into a single headline about whether LLMs beat humans. Our results suggest that this framing is too coarse. What matters is which information is available at forecast time and whether the model can access the signals that humans already aggregate. A more useful research direction is to characterize the division of labor between LLMs, retrieval, and human crowds rather than search for a single overall winner.

